\section{Exact Hitting-Time Potential and Deterministic Descent}
\label{sec:h-derandomization}

For a fixed satisfiable formula, let
\begin{equation}
H_F(x)=\mathbb E_x[T_{\mathrm{SAT}}].
\end{equation}
The Poisson equation implies
\begin{equation}
H_F(x)
=
1+\sum_yP_F(x,y)H_F(y)
\end{equation}
at every non-satisfying state.  Therefore at least one transition with positive
probability satisfies
\begin{equation}
\boxed{
H_F(y)\le H_F(x)-1.
}
\end{equation}

This gives an exact deterministic derandomization principle: if $H_F$ were
available in polynomial time, choose at each state an allowed successor having
minimum $H_F$.  The process must terminate after at most
$\lceil H_F(x_0)\rceil$ steps.

The complete finite-state hitting-time functions of the standardized
adversarial finalists were used to test this policy.

\begin{table}[t]
\centering
\caption{Deterministic descent guided by the exact hitting-time potential.}
\begin{tabular}{rrrr}
\toprule
$n$ & $\max H_F$ & Maximum deterministic steps & Steps/$n^2$ \\
\midrule
8 & 167.741 & 8 & 0.125 \\
10 & 424.181 & 10 & 0.100 \\
12 & 386.410 & 12 & 0.083 \\
14 & 340.104 & 14 & 0.071 \\
16 & 703.494 & 17 & 0.066 \\
\bottomrule
\end{tabular}
\end{table}

No violation of the guaranteed one-unit $H_F$ descent was observed; this
identity follows analytically from the Poisson equation.  The obstacle is
computational rather than existential: the present exact representation of
$H_F$ has $2^n$ states.  To obtain a deterministic polynomial-time SAT
algorithm, one needs a polynomially computable surrogate or representation
that preserves a comparable descent guarantee.
